\begin{decisione}{D0010}{Connessioni al database}{Accettata}{27/10/2026}

\contesto
PostgreSQL accetta un numero massimo di connessioni contemporanee (per impostazione predefinita 100, da verificare sulla versione usata). Ogni pod che scrive nel database ne tiene aperte alcune (\emph{pool}). Se il totale supera il limite, i pod nuovi ricevono errori di connessione e il peggioramento verrebbe attribuito per errore all'architettura. In V1, V2 e V4 scrivono le API; in V3 e V5 scrivono solo i worker, perché le API pubblicano sulla coda e non si collegano al database. Il caso peggiore è V5, con fino a 10 worker.

\scelta
\begin{itemize}
  \item \textbf{Regola:} il numero massimo di pod che si collegano al database, per le connessioni di ciascuno, non supera 80 (l'80\% del limite): il margine serve all'operatore del database, ai controlli di fine run e a PostgreSQL stesso.
  \item \textbf{Dimensione del pool:} 5 connessioni per pod, uguale per API e worker e per tutte le versioni, così il confronto è equo. Nel caso peggiore (V5, 10 worker) sono 50 connessioni.
  \item \textbf{Attesa di una connessione libera:} rientra nella deadline di 1,5~s di D0008. Se il pool è occupato e il tempo scade, l'API risponde 503.
  \item \textbf{Controllo prima dell'uso:} il pool verifica che la connessione sia viva e ricrea quelle morte, così dopo il riavvio del database il ripristino non dipende da connessioni vecchie.
  \item Il limite di PostgreSQL resta al valore predefinito (100) in tutte le versioni, anche con la replica di V4. Se la misura di capacità di F4 mostrasse che 5 connessioni per pod non bastano, si alza \texttt{max\_connections} per tutte le versioni insieme, annotandolo nella tabella dei parametri, invece di cambiare il pool in una sola.
\end{itemize}

\conseguenze
Con 5 connessioni per pod e 10 worker, V5 usa 50 connessioni su 100: il caso peggiore non esaurisce il database. Un pool di 10 porterebbe V5 a 100 e sarebbe al limite. Un pool piccolo può limitare il numero di richieste in corso verso il database, e durante F4 e F5 questo si vede come richieste respinte con 503 entro la deadline; è un effetto voluto, uguale per tutte le versioni.

\end{decisione}
